\section{Conclusione e Discussione Finale}
\label{sec:conclusion}

% per bayesiano: Nei paesi in forte espansione come Cina, Gran Bretagna,
%Australia, il PIL agisce da \textit{architetto} del successo olimpico: la sua
%crescita strutturale abilita investimenti che si manifestano nel medagliere con
%un anticipo di uno o due cicli olimpici, mentre l'assegnazione dei Giochi funge
%da catalizzatore di un processo già in corso. Dove invece le risorse esistono
%ma la loro allocazione settoriale è inefficiente, come nel caso italiano, il
%meccanismo si inceppa. Il successo olimpico è dunque il risultato di una
%maturazione sistemica che si sviluppa su orizzonti decennali: il picco
%casalingo non è un'anomalia, ma la punta visibile dell'iceberg
%dell'investimento sommerso.

Un fenomeno così complesso varia molto con ogni singola variabile, ma i
molteplici metodi impiegati evidenziano alcuni pattern.\\
Consolidata l'esistenza dell'\textit{host effect} (Sign-Test), la regressione
non lo esclude, ma lo decrementa considerando fattori storici e macroeconomici.
L'influsso del PIL viene anche individuato da Spearman (insieme alla
popolazione) e Changepoint Detection; quest'ultimo evidenzia una correlazione
compatibile l'anticipo di 7 anni per la candidatura, ma solo in casi in cui la
ricchezza sia investita efficacemente nello sport: ospitare spesso acuisce solo
un trend già esistente. Infine, la regressione evidenzia la non trascurabilità
dei fattori aggiuntivi e suggerisce come soprattutto quelli geografici fossero
più importanti in passato.